\section{Introduction}
\label{sec:introduction}

A language-model agent may invoke different models while solving a single task. A small model may handle routine actions, a larger model may respond to a failed tool call, and a subsequent decision may return execution to the smaller model. Evaluating such a strategy requires more than measuring the accuracy of each model on a fixed collection of prompts. A model choice can change the generated action, the environment state, the information available at the next decision, and whether another decision occurs. The scientific question is therefore sequential: for a specified population of tasks and an execution environment, what would happen if the agent followed a different rule for choosing models throughout execution?

This question arises in an active routing literature. Preference-based query routing and causal learning from observational query feedback address model selection for individual requests \citep{ong2025routellm,tsiourvas2025causal}. Multi-round routing, budget-aware agent routing, and routing embedded within an agent harness already address sequential model selection \citep{zhang2025routerr1,zhang2026baar,liu2026harness}. The Replay Gap directly motivates evaluation of the downstream consequences of changing a model during execution \citep{gonuguntla2026replay}. These studies establish that routing itself, including routing within an agent, is not a new problem. They motivate a complementary question: under what conditions can previously collected agent trajectories support a statistically interpretable comparison between routing policies?

A retrospective comparison between tasks that switched models and tasks that did not generally answers a different question. Failures can trigger escalation, so the switched group may contain harder tasks or histories already affected by earlier model choices. Conditioning only on the final switching indicator does not account for this feedback. Similarly, replacing a model label while preserving the recorded downstream observations can hold fixed variables that the intervention would change. A valid evaluation must specify both the routing decisions being changed and the subsequent execution mechanisms allowed to respond. This requirement also clarifies when replay is appropriate: a restored environment followed by an actual new continuation is an experiment, while an unchanged recorded future requires a separate invariance justification.

Dynamic treatment regimes and off-policy evaluation provide the relevant mathematical foundation. Longitudinal causal inference formalizes interventions in which later treatment choices depend on variables affected by earlier treatments \citep{robins1986new,murphy2003optimal}. Sequentially randomized designs make the assignment mechanism explicit \citep{murphy2005smart,klasnja2015mrt}. History-based reinforcement-learning estimators provide closely related value-identification and doubly robust estimation results \citep{jiang2016dr,thomas2016data,kallus2020drl}. These frameworks share the essential identification problem; adopting the language of treatment regimes does not confer an ability to handle feedback that off-policy evaluation lacks. Our formulation uses both perspectives to connect the deployed intervention, the data-collection design, and the scope of an inferential claim.

We define a \emph{dynamic agent regime} as a stochastic rule for selecting a versioned model configuration at prospectively specified eligible decisions. The configuration and execution contract include the model weights, quantization, prompt and context processing, decoding, tool interface, and relevant environment behavior. For a policy $\pi$, its value is $V(\pi)=E[G^{\pi}]$, where $G^{\pi}$ is a prespecified bounded return from a genuine execution under that policy. Terminal task success is one example; a declared utility incorporating measured resource use is another. The task distribution and resource horizon are part of the estimand. A claim about one benchmark or model version does not automatically extend to another.

This intervention definition permits model selection to be treated as a macro action. When the conditional execution kernels are unchanged, the likelihood ratio between target-policy and logging-policy trajectories contains the target-to-logging routing probabilities; the intervening generation and environment factors cancel. The analysis consequently does not require reconstructing token-level likelihoods. It does require that the logged probabilities describe the decisions actually executed, including availability restrictions and fallbacks. We establish an exact reduction from the fine execution process to blocks separated by eligible decisions, retaining complete boundary histories. The reduction applies when eligibility is determined from information available before assignment and intervening execution mechanisms agree. Its statistical horizon counts eligible opportunities, including decisions that retain the current model, rather than only observed switches.

For this observation model, we give self-contained identification and inference results using established longitudinal methods. The analysis includes policy-ratio weighting, the fixed-policy influence function, the doubly robust remainder, and conditions for cross-fitted uncertainty quantification \citep{bang2005dr,kallus2020drl}. Known sequential randomization removes the need to estimate assignment probabilities but does not remove variance caused by poor overlap. We therefore consider supported stochastic changes around a reference router, while distinguishing fixed-reference policies from interventions defined by an unknown behavior distribution \citep{kennedy2019incremental}. A finite collection of policies frozen before independent evaluation admits a conservative simultaneous improvement certificate under bounded-score conditions. This certificate concerns average policy value within the stated population and supported policy class.

We then address two design issues that arise when validating agent policies. First, restored branches can provide paired continuation outcomes at selected histories, but their sampling distribution need not match the histories encountered under a target policy. We derive a target-prefix-weighted augmented branch score, its variance, and a cost-constrained allocation rule within a specified class of independently selected branch experiments. The allocation has the classical variance-and-cost structure associated with stratified sampling \citep{neyman1934representative}; optimal data collection for policy evaluation is itself an established research topic \citep{li2023allocation}. The contribution here is the explicit connection between prefix occupancy, selective continuation measurement, and the value contrast being validated. It does not identify a full policy value from arbitrary branches.

Second, a policy may be evaluated under one execution contract and deployed under another. We provide a finite-horizon sensitivity bound for bounded complete-trace returns when the initial distribution and routing policy agree and stagewise execution kernels differ by declared uniform total-variation bounds. This is a full-history adaptation of simulation-lemma and coupling arguments \citep{kearns2002near,lobel2024simulation}. It quantifies the implication of specified kernel discrepancies; it does not infer those discrepancies from limited logs or identify the value of an unsupported model update.

The paper's contribution is thus an explicitly scoped evaluation design, with mathematical results connecting
prospective routing opportunities, selection-aware branch validation and execution-contract sensitivity.
The general identification and influence-function machinery is inherited; the branch allocation and drift
arguments specialize established principles. A descriptive open-weight coding study exposes a learned fixed
repair schedule, restricted decision opportunities and an unresolved offline/live discrepancy. It does not
demonstrate history-adaptive improvement, estimator superiority or calibrated uncertainty. The theoretical
analysis, this case study and a prospective validation protocol together identify what remains necessary to
determine whether the design provides useful precision or a resource advantage over direct execution.
